\chapter{Percurso Metodológico}

\section{Natureza, abordagem e estratégia da pesquisa}

A pesquisa será de natureza aplicada, com abordagem predominantemente qualitativa e apoio de dados quantitativos descritivos. Será desenvolvida como intervenção pedagógica de grupo único, com produções individuais em diferentes momentos da sequência. O desenho permite examinar mudanças, permanências e dificuldades nas explicações dos estudantes, mas não permite separar um efeito próprio da plataforma da mediação docente, das atividades e das interações em sala.

A estratégia principal será uma intervenção pedagógica. \citeonline{damiani2013discutindo} caracterizam esse tipo de investigação pelo planejamento e pela implementação de mudanças ou inovações pedagógicas, acompanhados da análise de seus efeitos. Essa formulação é compatível com o projeto porque os recursos da intervenção serão planejados para uma situação de ensino, aplicados no campo e examinados a partir das evidências produzidas.

A intervenção será compreendida como uma unidade formada pela sequência didática, pela plataforma, pela mediação docente e pelos instrumentos de produção de dados. Separar artificialmente esses elementos seria inadequado, pois, nesta pesquisa, a plataforma adquire o sentido previsto em conjunto com a pergunta proposta, a previsão solicitada e a discussão posterior.

Reconhecer essa unidade não autoriza atribuir qualquer mudança observada a um único componente. As produções dos estudantes serão lidas como evidências construídas em um percurso de ensino situado, no qual tempo de aula, infraestrutura, retomadas conceituais e interação entre duplas podem interferir na forma como uma explicação aparece. A posição do professor-pesquisador exigirá registro claro das condições de aplicação e prudência na interpretação dos dados.

A definição metodológica final deverá ser discutida com o orientador antes da submissão do projeto e do início da produção de dados. Essa decisão dependerá também das condições confirmadas do campo e das exigências éticas e institucionais aplicáveis.

\section{Contexto e participantes}

O campo previsto é uma turma de Ensino Médio de escola pública na qual o pesquisador atue como professor. A possibilidade inicial é uma turma vinculada à disciplina Projeto Permanente por Afinidade, em escola de tempo integral. Esse cenário não está garantido para o próximo ano letivo, pois a lotação docente e a composição das turmas poderão mudar.

Se necessário, a sequência poderá ser aplicada em outra turma do Ensino Médio sob docência do pesquisador, inclusive em um componente de Matemática, desde que sejam preservadas as condições pedagógicas e éticas definidas para a pesquisa.

Os participantes serão estudantes iniciantes em \textit{machine learning}. Não haverá exigência de programação prévia. A sequência pressupõe leitura de tabelas e gráficos elementares, localização de pontos em eixos cartesianos, comparação de valores e mobilização de noções iniciais de média, distância, função e probabilidade quando esses conhecimentos forem retomados durante os módulos.

O número de participantes, a série, o componente curricular, a escola e o calendário de aplicação serão informados após a confirmação do campo.

\section{Etapas da intervenção}

O primeiro momento corresponderá ao estudo das referências e ao aperfeiçoamento da plataforma, dos exemplos sintéticos e dos instrumentos. Em seguida, a sequência didática será preparada de acordo com a turma e as condições de infraestrutura disponíveis.

Prevê-se que os estudantes trabalhem em duplas, cada uma com um computador do laboratório de informática. No campo inicialmente previsto, o laboratório dispõe de 20 máquinas; em outro campo, será necessário verificar as condições disponíveis. O trabalho em duplas cria espaço para que os estudantes comparem hipóteses diante do mesmo experimento. As previsões e explicações selecionadas para análise serão registradas individualmente antes e depois da exploração em três encontros. Os registros docentes informarão quando uma discussão em dupla, uma retomada conceitual ou uma condição de infraestrutura puder ajudar a interpretar essas produções. Os oito encontros correspondem a uma duração viável para esse campo e exigirão atenção às retomadas e ao ritmo efetivo da turma.

Ao final de cada encontro, o professor-pesquisador organizará uma síntese breve que recupere os elementos discutidos pela turma.

As sínteses de encerramento não terão a função de produzir uma resposta correta única para ser copiada pelos estudantes. Elas servirão para recuperar o que foi observado, explicitar relações ainda confusas e preparar a passagem para o encontro seguinte. Caso a turma confunda uma probabilidade com garantia de ocorrência, por exemplo, a retomada poderá recuperar o papel do limiar e distinguir estimativa de certeza antes de avançar para outro modelo.

O registro docente indicará que a retomada ocorreu e quais dificuldades a motivaram, sem substituir as respostas individuais como fonte de evidência da compreensão. Essa separação ajuda a preservar o vínculo entre a mediação necessária ao ensino e a análise cuidadosa do que cada participante conseguiu elaborar.

A plataforma não registrará respostas ou dados pessoais. As respostas usadas na pesquisa serão produzidas nos instrumentos descritos a seguir.

\section{Instrumentos de produção de dados}

Antes da sequência, os participantes realizarão um diagnóstico individual breve. O instrumento terá um núcleo de quatro ou cinco questões sobre dados e finalidade da tarefa, consequência de uma alteração, interpretação de representação e limite de uma decisão automatizada. Ele servirá para registrar o ponto de partida das explicações e não comporá a avaliação escolar.

Nos encontros 2, 4 e 6, cada participante preencherá um registro curto com duas solicitações: indicar o que prevê antes de alterar o controle e explicar o resultado depois da exploração. A escolha desses encontros permite acompanhar uma tarefa de previsão numérica, uma tarefa de classificação e uma tarefa de agrupamento, sem exigir produção escrita em todos os módulos.

O instrumento final individual terá dois blocos. O primeiro retomará o núcleo do diagnóstico por meio de questões equivalentes em operação cognitiva, extensão e dificuldade. Os contextos poderão mudar, mas cada par de questões deverá solicitar a mesma relação conceitual. O segundo bloco será um desafio de transferência. Nele, o estudante analisará uma situação diferente das usadas nas aulas, identificará os dados necessários, reconhecerá o tipo de tarefa, indicará um modelo compatível e explicará uma possível consequência da alteração de um dado, parâmetro ou regra.

O desafio de transferência não será tratado como equivalente ao diagnóstico. Sua função será examinar a mobilização dos conceitos em outro contexto. Um questionário breve de percepção poderá registrar clareza, dificuldade e utilidade atribuída às atividades e à plataforma. Essas respostas serão analisadas como apreciação do produto, e não como demonstração de aprendizagem.

Depois de cada aula, o professor-pesquisador produzirá um registro em tópicos sobre dificuldades recorrentes, intervenções realizadas, tempo empregado e condições de uso. As perguntas inseridas no site terão função formativa imediata. Como suas respostas não permanecem salvas, não serão usadas diretamente como dados da pesquisa.

Antes da aplicação, uma matriz relacionará cada questão comparável ao objetivo de aprendizagem, à atividade em que o conteúdo foi trabalhado, à dimensão da rubrica e à questão correspondente no outro momento. Os instrumentos serão submetidos à leitura crítica do orientador e de um professor com experiência no nível de ensino. Quando for viável, a clareza dos enunciados será verificada com pessoa de perfil próximo ao público que não participe da pesquisa.

Os instrumentos também precisarão ser compatíveis com o tempo escolar disponível e com o fato de que os participantes serão estudantes do Ensino Médio em uma atividade inserida no cotidiano da turma. Não se pretende aplicar teste extenso que interrompa a sequência ou transforme a intervenção em rotina de avaliação formal.

Perguntas abertas curtas, acompanhadas de gráfico, tabela ou situação-problema, podem oferecer indícios mais úteis que grande quantidade de itens objetivos, desde que o enunciado deixe claro o que se pede. Antes da aplicação, o professor-pesquisador verificará se cada questão está ligada a um objetivo de aprendizagem e a uma atividade efetivamente realizada.

Caso algum conceito não tenha sido abordado nas condições previstas, ele não deverá ser tratado como ausência de compreensão no instrumento final. Essa precaução mantém os instrumentos vinculados ao percurso de ensino e protege a análise contra interpretações que ultrapassem as evidências disponíveis.

\section{Procedimentos de análise}

As produções serão organizadas por códigos, sem identificação nominal dos participantes. A rubrica terá três níveis, de 0 a 2, em quatro dimensões: dados e finalidade; alteração e consequência; representação e conceito; possibilidades e limites. O nível 0 reunirá respostas ausentes, incompatíveis ou sem relação justificável com a situação. O nível 1 indicará relação parcial entre os elementos. O nível 2 indicará uma explicação consistente com o nível introdutório definido para a sequência. Cada classificação será acompanhada por uma justificativa curta.

A primeira etapa consistirá na leitura integral das respostas e na revisão dos exemplos que delimitam os três níveis. Em seguida, será organizada uma matriz por participante, momento e dimensão. Os códigos apoiarão a organização do corpus e não serão apresentados aos estudantes como nota.

Uma resposta poderá apenas indicar que ``a classificação mudou'', enquanto outra relacionará a mudança ao limiar escolhido ou à proximidade entre exemplos. A diferença entre esses modos de explicar será mais informativa que uma oposição simples entre certo e errado.

As questões comparáveis do diagnóstico e do instrumento final serão colocadas em relação. A análise registrará permanência no mesmo nível, avanço ou redução, sempre com leitura da justificativa escrita. Uma tabela de transição por dimensão mostrará quantos participantes seguiram cada percurso. Os registros intermediários ajudarão a descrever como certas relações apareceram durante as aulas. O desafio de transferência será analisado separadamente, segundo as mesmas dimensões quando forem aplicáveis.

Frequências absolutas e relativas serão apresentadas com seus denominadores. Respostas ausentes, instrumentos incompletos e participantes sem par entre os dois momentos serão informados, sem substituição de valores. Essas medidas descreverão o conjunto e não sustentarão inferências estatísticas ou generalizações para outras turmas.

O texto final reunirá sínteses por dimensão, tabelas de transição, frequências descritivas e poucos excertos anonimizados que representem modos distintos de explicação. Os registros docentes serão usados quando uma dificuldade recorrente, uma retomada matemática ou uma condição de infraestrutura ajudar a interpretar as produções escritas. O questionário de percepção será apresentado em seção própria e não será combinado com os níveis da rubrica.

A interpretação considerará as condições em que a atividade ocorreu, mas não usará registros do professor para completar uma resposta individual. Quando for viável, um segundo avaliador classificará uma parcela das produções. As divergências serão discutidas com base nas definições e nos exemplos da rubrica, e os ajustes de critério serão registrados antes da análise final.

\section{Aspectos éticos}

A produção de dados somente poderá começar após a confirmação das exigências institucionais e éticas aplicáveis ao projeto. A Resolução CNS nº 510/2016 dispõe sobre pesquisas em Ciências Humanas e Sociais que utilizam dados obtidos diretamente de participantes ou informações identificáveis, enquanto a Resolução CNS nº 466/2012 estabelece diretrizes gerais para pesquisas envolvendo seres humanos \cite{brasil2016cns510,brasil2012cns466}.

Como a intervenção poderá envolver adolescentes, os procedimentos éticos deverão ser definidos conforme a orientação institucional e a tramitação exigida antes da aplicação.

As respostas serão anonimizadas por códigos, e documentos com informações de identificação, quando necessários, serão mantidos separadamente dos instrumentos de análise. Dados brutos não serão inseridos no código da plataforma nem publicados em repositório aberto. A participação será voluntária. Estudantes que não participarem da produção de dados deverão ter acesso a atividade pedagógica equivalente, sem prejuízo da experiência escolar.
